\section{Scientific Supplement}
\label{sec:supplement}

This supplement gives the statistical definitions, extended comparisons, original LOCAL pairs, eligibility rules and numerical evidence scope. The \href{https://github.com/leecdiang/HeteroSpec/blob/a634e552366dec9b95ad1003dabdf1028a8ee63d/papers/taco/20261118_atelier_integration/artifact/SCIENTIFIC_DATA_INDEX.md}{scientific data companion} indexes every measured request and paired round, including all 864 RTX confirmation requests, 96 A5000 confirmation observations and earlier paired records. It also provides the source identities and CPU reconstruction instructions.

\FloatBarrier
\subsection{Original selected transactions and dependent successors}
\label{supp:local}
The original LOCAL study uses CodeContests/train/161 histories beginning at 137 and 253, with two paired repeats each. Successors begin at 139 and 255 after their own preceding commits. The verification-through-commit clock includes owned-row copying, decisions, trim and synchronization; history reconstruction, loading, live drafting and payload persistence are outside it. Figure~\ref{fig:local} and Table~\ref{tab:local-pairs} retain every clean pair.

\paperfigure{f03_transaction}{Original LOCAL measurements for two fixed histories and two paired repeats (four first-transaction pairs), with C4 on, P4 off and frozen proposals. Open squares are SIMPLE; filled circles are LOCAL; segments join matched arms. (a) First-transaction verification-through-commit time. (b) Separately clocked projection and attention event time for each history; both panels use milliseconds. Device-event times are not added to the transaction clock. Four dependent successor pairs are retained in Table~\ref{tab:local-pairs}.}{The upper panel shows four first-transaction pairs: two crossing repeats at base 253 and two noncrossing repeats at base 137. Open SIMPLE squares and filled LOCAL circles are connected within each pair, with milliseconds and LOCAL/SIMPLE ratios printed at right. The lower panel separately compares projection and attention event times for both histories. These device-event costs use their own clock; they are not components to add to the transaction wall clock.}{fig:local}

\papertable{s01_local_pairs}{Every original LOCAL pair, including successors from each arm's own committed state.}{tab:local-pairs}

The crossing first-transaction geometric ratio is 0.77395294 (22.96\% and 22.24\% reductions in its two repeats); its successor ratio is 0.79822856. The noncrossing first ratio is 0.98022489, combining a 1.25\% increase and a 5.11\% reduction; its successor ratio is 0.99980702. These are descriptive fixed-history estimates without a population interval. Successors test state dependence and do not enlarge the independent sample count. Table~\ref{tab:diagnostics} reports the separate device-event diagnostics.

\papertable{s02_diagnostics}{Independent device-event diagnostics for the two first-transaction histories. Times and counts use the profiler's own clock.}{tab:diagnostics}

At the crossing, SIMPLE executes 449 q3 and 385 q1 events in each preparation/dot category; LOCAL executes 449 q4 events. The q1 path partly uses gate/up fusion. LOCAL adds 64 concatenations, totaling 0.134\,ms; the exported memcpy aggregate rises from 0.361 to 0.494\,ms. Profiler/cold-call costs can enter these diagnostics. The full lowering changes widths, fusion, launch organization, views and joins together.

\FloatBarrier
\subsection{A5000 confirmation: time blocks, endpoints and state}
\label{supp:confirmation}
Panel A repeats the selected 32B histories under the original numerical build, C1024/F16, C4 on, P4 off and frozen proposals. Noncrossing is primary; crossing and paired interaction are supporting (Table~\ref{tab:confirmation-effects}). Panel B uses the original 8B Q8\_0 target-only two-turn consumer; its primary endpoint gives equal log weight to the two fixed conversations. Each panel has 12 paired time blocks, analyzed separately from the original study.

\papertable{s14_confirmation_effects}{All six A5000 confirmation endpoints and their prespecified roles. The interaction is a ratio of ratios.}{tab:confirmation-effects}

For Panel A, let $\ell_{hb}$ be the paired log time ratio for history $h$ in block $b$. The direct estimate is $\exp(12^{-1}\sum_b\ell_{hb})$, and the interaction is $\exp(12^{-1}\sum_b[\ell_{\mathrm{cross},b}-\ell_{\mathrm{noncross},b}])$. Panel B uses $\exp(12^{-1}\sum_b[\ell_{0b}+\ell_{1b}]/2)$ overall and the corresponding single-conversation means for supporting endpoints. The 10,000 bootstrap draws (seed 20261117) resample whole blocks jointly, preserving paired cases; individual 95\% percentile intervals use linear interpolation.

Adjacent block starts are at least 1,978 seconds apart. The campaign crosses midnight in \mbox{Asia/Shanghai}, with eleven blocks before midnight and one after it. Intervals describe the fixed cases in this observed window. First/last-half, order, date and lag-one summaries remain sensitivity descriptions. Tables~\ref{tab:confirmation-a-pairs} and~\ref{tab:confirmation-b-pairs} retain every block.


All twelve crossing pairs favor LOCAL; five noncrossing pairs favor SIMPLE. In B, seven C0 pairs, six C1 pairs and five equal-conversation block aggregates favor SIMPLE. The equal-conversation log estimate differs from the older natural panel's ratio of totals.

A's crossing changes SIMPLE q3+q1 into LOCAL q4 with attention 3+1; the successor's 1+3 partition triggers the full-vector guard. The noncrossing successor reuses the GGML graph. CUDA graphs are disabled, independently of this graph reuse. Crossing mean first-transaction time changes by $-21.555$\,ms: $-21.731$\,ms decode, $+0.176$\,ms other host work, and $+0.000513$\,ms trim/synchronization. Its dependent continuation separately saves 20.858\,ms. Prefix preparation adds 4.924\,ms and the scientific wrapper adds 62.252\,ms outside the transaction clock; the wrapper has unequal-depth checks and is not an application endpoint. No new separated projection/attention profiler clock was collected.

\papertable[!htbp]{s15_confirmation_A}{Every 32B confirmation pair: first-transaction verification-through-commit milliseconds. Noncrossing is primary and crossing supporting.}{tab:confirmation-a-pairs}

\begin{samepage}
B preserves the parser, wire/client history and re-tokenization chain, C4096, two freely generated greedy turns and cap128 per turn. Full prompts are 92/250 tokens for C0 and 27/551 for C1. Second turns reuse 88/23 tokens, leaving 162/528 new prompt tokens; every turn emits 128 tokens. C0 retains 263 physical calls and 508 target rows in either arm. C1 has 273 SIMPLE versus 272 LOCAL calls at 809 rows; LOCAL starts 247 and 503 span internal $9+23$ boundaries. These widths belong to this original Q8 natural path. Table~\ref{tab:confirmation-costs} gives the response-plus-reset phase differences.\par
\end{samepage}

\papertable[!htbp]{s16_confirmation_B}{Every 8B confirmation pair: complete two-turn response-plus-reset seconds. C0/C1 identify conversations in this panel.}{tab:confirmation-b-pairs}

\papertable{s17_confirmation_costs}{Arithmetic mean LOCAL-minus-SIMPLE phase differences in milliseconds within the registered response-plus-reset clock.}{tab:confirmation-costs}

The 212 requests comprise 12 qualification, 96 fresh same-arm warmups, 96 formal observations and eight diagnostics. A and B each contain 48 formal arm observations and 106 total units; a B unit is a two-turn episode. Technical replacements were disabled, and there were no failures, interruptions or reruns. Recorded work is 91,319 target rows and zero draft rows. Original-machine captures check 954 A owned/frontier vectors and 230 canonical checkpoints, plus 27,136 B vocabulary vectors and 444 active-KV snapshots. Capture depth varies by role; these are recorded counts, not per-request extrapolations. Public numerical access is specified in Table~\ref{tab:access}.

\FloatBarrier
\subsection{RTX confirmation after shared activation preparation}
\label{supp:rtx-confirmation}

The six-block RTX confirmation retains the earlier bridge's numerical core, native workers and models. Its 24 exposed roots have 12 rounds each, three arms and two rounds per root/block. Every root receives each of the six arm orders twice; every domain/block/arm/position cell contains eight observations. Each root-round uses three fresh workers with same-arm warmups and measured requests. The observed campaign runs from 17 November 2026 21:33 to 18 November 06:56 in \mbox{Asia/Shanghai}, with at least 1,805.7 seconds between successive blocks. There are 24 distinct inputs and six temporal clusters.

C1/B is the sole prespecified primary comparison; C4/C1 and C4/B are descriptive secondary contrasts (Table~\ref{tab:rtx269-effects}). For root $i$, block $b$ and its two rounds $j$,
\begin{equation}
 \widehat r_{\mathrm{RTX}}=\exp\!\left[\frac{1}{24\cdot6}\sum_{i=1}^{24}\sum_{b=1}^{6}\left(\frac12\sum_{j=1}^{2}\log\frac{T^{\mathrm{candidate}}_{ibj}}{T^{\mathrm{control}}_{ibj}}\right)\right].
 \label{eq:rtx-estimator}
\end{equation}
The 10,000 prespecified draws independently resample roots within balanced code/math strata and whole time blocks, using common draws across all contrasts and linear percentile endpoints. The primary interval spans a 0.48\% reduction through a 2.19\% increase and is not contained in the prespecified $\pm2\%$ practical band; the data support neither a resolved benefit direction nor equivalence within that band.

\papertable{s18_rtx269_effects}{RTX confirmation primary and descriptive complete-request estimates, with crossed root/time-block intervals.}{tab:rtx269-effects}

The conditional root-only C1/B interval is [1.001432, 1.015407], holding the six blocks fixed. Tables~\ref{tab:rtx269-blocks} and~\ref{tab:rtx269-roots} retain every block and root. One C1/B block favors C1 and five favor B; all 24 root aggregates and all six block aggregates favor C4 against either scalar control.

\papertable{s19_rtx269_blocks}{All six RTX time blocks in execution order: within-block geometric time ratios.}{tab:rtx269-blocks}
\papertable{s20_rtx269_roots}{All 24 RTX roots: per-root geometric time ratios over 12 rounds.}{tab:rtx269-roots}

Table~\ref{tab:rtx269-phases} closes the arithmetic phase account to complete service, retaining prefill and drafting. C1/B saves 100.43\,ms in verification, adds 173.57\,ms in prefill and 7.43\,ms in draft work, and changes remaining stages by $-0.06$\,ms: $+80.52$\,ms net. C4/C1 saves 1,465.29\,ms in verification while prefill adds 63.21\,ms, for a 1,411.60\,ms net request saving. These are arithmetic differences, separately defined from geometric effects.

\papertable{s21_rtx269_phases}{RTX confirmation arithmetic phase differences, candidate minus control, within the complete resident-request clock.}{tab:rtx269-phases}

C4 records 2,273,400 verification admissions and 2,652 prefill admissions. All prefill admissions occur at the q3 tail of the 163-token MATH/counting\_and\_probability/train/398 prompt, in its 12 measured requests. B and C1 have zero few-column admissions. The treatment includes both target verification and that observed prefill exposure. Prompts span 140--757 tokens; visible outputs 42--128; decision sequences including EOG 43--128. The ledger contains 684 cap-terminated and 180 EOG requests; separate initial-EOG and complete logical q1-tail qualification remain outside this panel.

The ledger has 1,764 requests: 864 formal measured, 864 formal warmups, 24 eligible qualification and 12 resource-excluded requests. The excluded units are one complete six-request qualification unit and one complete six-request formal root-round unit, replaced under the frozen technical allowance; all original records remain. Eligible qualification contains twelve capture warmups and twelve clean requests, excluded from estimation even where an internal field says measured.

\begin{samepage}
Original-machine qualification compares fresh candidate captures with stored pristine references on two sentinel roots. Its overlapping counts are 1,932 pristine-target vector comparisons, 3,768 target/draft vector comparisons and 456 canonical checkpoints. Every timed request validates complete-state/frontier identities, pending and transaction state, actual target/draft work, dispatch and graph/view identity. Public material includes complete final frontiers, native records and bindings; large full-vector and KV payloads remain on the original host with manifests.
\par
\end{samepage}

\FloatBarrier
\subsection{Complete-request effects and phase accounts}
\label{supp:effects}
The earlier modern request studies use Eq.~\eqref{eq:estimator}: paired log ratios are averaged over a root's rounds before giving roots equal weight. Whole-root bootstrap retains rounds within code/math strata. The 10,000-draw percentile intervals use linear interpolation, seed 20261117 for A5000 C4 and 20261116 for the other modern panels. Table~\ref{tab:direct-effects} includes every direct contrast from the five principal application studies.

\papertable{s11_direct_effects}{All direct modern request effects, retaining each independent study and its contemporaneous control.}{tab:direct-effects}

\paperfigure{f08_c4_initial}{Initial C4 studies with SIMPLE fixed and P4 off, separate from the 12-round RTX confirmation. (a) Twelve A5000 roots: paired C1/C4 geometric complete-request seconds over two rounds. (b) Twenty-four initial RTX roots: geometric time ratios for C1/B, C4/B and C4/C1 over two rounds, with equal time at one. Summaries give equal-root estimates and individual 95\% stratified-root intervals. The two studies are unpaired.}{The upper panel uses seconds and links the C1 open square to the C4 filled circle for each of 12 A5000 roots. The lower panel uses time ratios for 24 initial RTX roots in each of three direct comparisons. C1/B has an estimate of 1.008 and interval 0.994 to 1.022, crossing equal time; both C4 comparisons lie below one. Each initial study has two rounds and a separate control and interval; neither is pooled with the later RTX confirmation.}{fig:c4-initial}

Table~\ref{tab:phase-costs} reports nonoverlapping arithmetic mean phase differences, including residual host cost, that close to each named service boundary. These differences and geometric effects have separate estimands. CUDA resident clocks exclude model/context setup and post-clock exports; B570 includes normal teardown.

\papertable{s03_phase_costs}{Signed arithmetic mean phase differences in milliseconds. Each column is one same-study SD/SD comparison.}{tab:phase-costs}

A5000 has 232 P4-entering prefill calls across TO\_P4 and SD\_P4, all aligned q32 within one physical extent. Each RTX P4 arm records 140,114 q32, 442 q8 and 2,652 other-width projection dispatches, with zero q16 in the recorded buckets: 143,208 dispatches per arm. Dispatch counts and outer-call counts are distinct units; these records locate actual workload exposure.

\FloatBarrier
\subsection{Natural service and the complete M1 design}
\label{supp:natural}
M1 crosses WHOLE/LOCAL with loop4/packed projection organization, retaining four sessions within each of 16 roots. Every arm has 2,224 logical cycles. WHOLE executes 2,244 physical verification calls: 2,204 q4, 16 q2, 12 q1 and 12 q3. LOCAL executes 2,224 q4 calls, removing 20 extra calls (0.8993\% of cycles).

\papertable[!htbp]{s06_m1_family}{M1's five-comparison family. Direct contrasts use ratios of totals; interaction uses root-geometric ratios.}{tab:m1-family}

For a direct effect, time is summed across each root's four sessions before paired-root bootstrap and the ratio of totals. The interaction is
\[
 \exp\!\left[\frac{1}{16}\sum_i\log
 \frac{L_{P,i}/W_{P,i}}{L_{L4,i}/W_{L4,i}}\right].
\]
The bootstrap uses 10,000 paired-root draws (seed 19720261107) and 99\% marginal intervals for the five-comparison Bonferroni family. LOCAL is slower on five loop4 roots and ten packed roots; packing improves every root under either cut policy (Table~\ref{tab:m1-roots}). An interaction above one indicates attenuation of LOCAL's effect under packing.

\papertable{s07_m1_roots}{Every M1 paired root. Times sum four sessions in seconds; each row gives its direct cut-policy ratios.}{tab:m1-roots}
\papertable{s05_natural}{Other historical natural LOCAL/WHOLE panels with their original units, clocks and estimators.}{tab:natural}

The historical RTX panel has four roots and three sessions (12 paired records). The A5000 target-only panel has two conversations and two rounds (four episodes per arm); its controller clock includes reset and shared sampling. These descriptive panels retain their original estimates, separately from the selected 32B transactions and later confirmations.

\FloatBarrier
\subsection{Eligibility, technical failures and scientific stops}
\label{supp:ledger}
Table~\ref{tab:ledger} separates recorded roles before estimation. Qualification remains excluded even when an internal field says measured. The RTX P4 resource-marked block and its registered whole-block replacement are both recorded. Every eligible measured root enters the estimator.

\papertable{s04_ledger}{Request-role accounting. Timing-table scope and full-campaign completion are distinct.}{tab:ledger}

The A5000 C4 table has 102 requests plus 12 independent width-qualification requests: 114 combined. Its six sentinel captures, 48 warmups and 48 formal requests retain separate roles. A duplicate qualification-reader callback was corrected offline before formal reservation; captured requests were not rerun.

LOCAL records 27 reservations, 25 completions and two failed checkers. One checker closed observation counters; after observer repair the planned same-arm warmup supplied qualification. The other expected unfused q1 counts; its expectation was corrected against recorded fusion. A missing profiler export used one registered diagnostic replacement. All eight formal measured requests remained unchanged. The B570 formal timing table has 384 requests; 108 qualification requests and earlier technical attempts lie outside it. The full campaign has 514 launches, 513 completions and one failure. Its two-root pilot has eight measured, eight warm and 15 qualification/capacity requests. The predeclared complete-service criterion stopped the pilot, so the planned 24-root continuation did not run; this was a scientific stop, not a technical exclusion.

\FloatBarrier
\subsection{Numerical qualification and accessible saved objects}
\label{supp:access}
The A5000 C4 inventory contains 974 target and 724 draft vocabulary objects. Each candidate arm compares 484 conditional rows with a common reference, giving 968 candidate-to-reference comparisons. LOCAL contains 216 owned, 48 captured-native and 27 incoming-frontier vocabulary objects (291 total), plus 72 state objects. These original-machine inventories are broader than the public subset.

\papertable{s08_access}{Numerical evidence scope: original-machine comparisons and the independently readable saved-object subset.}{tab:access}

The frozen public archive includes q16 full KV at base32, active lengths 48/49 for R/C4/P4, and selected LOCAL crossing rows. The later release adds within-block LOCAL final KV and a retained 8B prefix. Preserved byte checks cover 32 owned-row, four incoming-frontier and four selected B-vector comparisons, plus five canonical state bindings across 24 unique payloads. B's frontier is invalid; its object is a retained prefix, not a completed service endpoint. Crossing LOCAL full KV and the remainder of the roughly 43\,GB numerical archive remain on the original host. Saved-byte checks compare captured objects without executing the model. Separate q12/q20/q28 qualification, complete logical q1-tail/EOG capture, and RTX q16/q24 qualification lie outside the exercised panels in Table~\ref{tab:qualification}.

\FloatBarrier
\subsection{The stopped B570 draft-placement pilot}
\label{supp:pilot}
The pilot moves 112 draft main matrices and its LM head to Vulkan0, retaining a CPU-mapped input embedding. Its two predeclared roots have SD\_C/TO\_C ratios 1.90760485 and 2.83651421, with geometric ratio 2.32614451. Mean draft prefill, synchronization and proposal cost is 14.807324\,s, exceeding the 1.642558\,s target decode-to-verification saving (Table~\ref{tab:pilot}). This cost account belongs to the two-root placement test; the formal 24-root CPU-draft study has its own 2.512087 service ratio.

\papertable[htbp]{s09_pilot}{The separate B570 draft-placement pilot: two roots, one round and a predeclared complete-service stop.}{tab:pilot}

\FloatBarrier
\subsection{Additional historical results and their scope}
\label{supp:implementation}
The historical register retains a 32B scalar-reference application that is 14.00\% slower than target-only across four roots even though its speculative-path change reduces time by 9.61\%. A separate shared-activation scalar S3 panel gives a 1.5071\% reduction [1.2436\%, 1.7554\%] across 12 roots and two observed dates. Its B/C paths and target-only clock differ from the modern C1/C4 bridge.

A B570 fixed-append study on SmolLM2-1.7B uses C1024/4096, q2/q4/q32 and three paired repeats per cell. Its $1+(q-1)$ crossings give 20.34--44.68\% reductions, while within-block LOCAL takes 0.43--3.74\% more time. Earlier Qwen8B/SmolLM2 q32 cells with $16+16$ crossings include slower LOCAL results (Table~\ref{tab:historical-geometry}). Each cell retains its capacity, prefix, requested outputs and clock.

\papertable[htbp]{s12_historical_geometry}{Earlier Vulkan geometry observations, including slower LOCAL cells and their original fixed-cell summaries.}{tab:historical-geometry}

Earlier Qwen8B, Qwen14B and Llama resident studies report 17.21\%, 25.54\% and 11.52\% reductions against their fastest measured target-only controls. They use composite speculative policies and their own family-adjusted intervals; two Llama roots are slower. Llama's short-worker total reverses the resident advantage, 368.822 versus 355.730\,s, after loading, checks and primers enter the lifetime. Another historical RTX q8 equal-work panel saves 3.94--11.91\% for the same eight known rows, while live k7 costs 13.58--14.59\% more than k3 at equal emitted-token counts but more verification rows and draft proposals.

\papertable[htbp]{s13_historical_applications}{Historical resident effects and the separate worker lifetime. Fastest resident and worker controls are selected within their own measured configurations.}{tab:historical-applications}

The historical studies use their original analyses. The artifact supplies their source identities, complete cell records, and the scope of saved-record reconstruction.
